\section{Реализация и тестирование программного комплекса}

В данной главе представлено детальное описание реализации ключевых модулей программного комплекса, а также приводятся и анализируются результаты, полученные в ходе его тестирования на целевой платформе (Apple MacBook Air с процессором M1).

\subsection{Реализация пользовательского интерфейса}

Пользовательский интерфейс (GUI) является точкой входа для взаимодействия пользователя с системой. Он спроектирован с упором на простоту, наглядность и отзывчивость. Основой для реализации послужила библиотека Fyne. Главное окно приложения состоит из контейнера с вкладками (\texttt{container.NewAppTabs}), который организует основной функционал по четырем направлениям: «Мониторинг», «Производительность», «Память» и «Многозадачность».

\subsubsection{Панель мониторинга}

Панель мониторинга (\texttt{DashboardPanel}) — наиболее динамичная часть интерфейса. Ее ключевая особенность — фоновая горутина, которая с интервалом в 500 миллисекунд опрашивает модуль \texttt{profiling} и асинхронно обновляет виджеты. Информация сгруппирована в карточки (\texttt{widget.NewCard}) для логического разделения. Особого внимания заслуживает реализация отображения загрузки ядер ЦП. При инициализации панели определяется количество ядер и их тип (P-core/E-core для macOS). Затем для каждого ядра создается свой прогресс-бар, который обновляется в цикле.

\subsubsection{Панели тестирования}

Все три панели с тестами (\texttt{BenchmarkPanel}, \texttt{MemoryPanel}, \texttt{ConcurrencyPanel}) имеют схожую структуру, основанную на универсальном обработчике тестов \texttt{TestRunner}. Логика их работы заключается в следующем:
\begin{enumerate}
    \item При создании панели формируется список тестов. Для каждого теста определяется его название, описание и функция, реализующая его логику.
    \item На панели размещаются кнопка для запуска всех тестов и общий прогресс-бар.
    \item При нажатии на кнопку запускается горутина, которая итерируется по списку тестов и для каждого из них вызывает \texttt{TestRunner}. Это гарантирует, что интерфейс не будет заблокирован на время выполнения длительных вычислений.
    \item \texttt{TestRunner} многократно выполняет переданную ему функцию теста, собирает статистику (min, max, avg) и сообщает о прогрессе через специальный канал. Данные о прогрессе агрегируются и отображаются на общем прогресс-баре.
    \item По завершении каждого теста его результат форматируется и отображается в соответствующей карточке на панели.
\end{enumerate}

\subsection{Реализация и результаты тестов}

\subsubsection{Тесты производительности ЦП}

\paragraph{Целочисленная арифметика.} Тест выполняет большое количество базовых операций (сложение, вычитание, умножение, деление) над 64-битными целочисленными переменными в цикле. Этот тест хорошо отражает базовую производительность процессора на задачах общего назначения, таких как работа с циклами, индексами массивов и т.д.

\begin{lstlisting}[language=Go, caption={Реализация теста целочисленной арифметики}]
const numOps = 10000000

func IntegerBenchmark() (float64, string) {
	start := time.Now()
	var a, b, c, d int64 = 100, 200, 300, 400
	for i := 0; i < numOps; i++ {
		a = b + c
		b = a - d
		c = a * b
		d = c / (b + 1) // +1 чтобы избежать деления на ноль
	}
	elapsed := time.Since(start)
	return float64(elapsed.Nanoseconds()) / numOps, "ns/op"
}
\end{lstlisting}

\textbf{Результат на Apple M1:} в среднем \textbf{0.25 ns/op}. Это означает, что процессор выполняет одну такую итерацию цикла примерно за 1-2 такта, что свидетельствует о высокой эффективности его конвейера.

\paragraph{Арифметика с плавающей запятой.} Аналогичный тест, но с использованием чисел с плавающей запятой двойной точности (\texttt{float64}). Производительность в этом тесте критически важна для научных вычислений, 3D-графики, обработки сигналов и машинного обучения.

\textbf{Результат на Apple M1:} в среднем \textbf{0.30 ns/op}. Результат очень близок к целочисленному, что говорит о высокой производительности блока FPU (Floating-Point Unit) в процессоре.

\subsubsection{Тесты иерархии памяти}

\paragraph{Задержка доступа.} Для этого теста используется алгоритм «погони за указателями» (pointer chasing). Его реализация заключается в создании связанного списка, узлы которого псевдослучайно перемешаны в большом массиве. Затем в цикле происходит обход этого списка. Так как адрес следующего узла становится известен только после чтения текущего, это полностью нивелирует эффект аппаратной предвыборки данных (prefetching) и позволяет измерить «чистую» задержку.

\begin{lstlisting}[language=Go, caption={Реализация теста задержки памяти}]
// node - узел связанного списка
type node struct {
	next *node
}

// chase - функция обхода списка
func chase(p *node, steps int) *node {
	for i := 0; i < steps; i++ {
		p = p.next
	}
	return p
}

// measure - основная измерительная функция
func measure(sizeBytes int, steps int, repeats int) float64 {
	size := sizeBytes / int(unsafe.Sizeof(node{}))
	results := make([]float64, repeats)
	for r := 0; r < repeats; r++ {
		// Создаем массив и перемешиваем его, чтобы сформировать
		// псевдослучайный связанный список
		nodes := make([]node, size)
		perm := rand.Perm(size)
		for i := 0; i < size-1; i++ {
			nodes[perm[i]].next = &nodes[perm[i+1]]
		}
		nodes[perm[size-1]].next = &nodes[perm[0]]
		head := &nodes[perm[0]]

		start := time.Now()
		_ = chase(head, steps) // Выполняем погоню за указателями
		elapsed := time.Since(start)
		results[r] = float64(elapsed.Nanoseconds()) / float64(steps)
	}
	sort.Float64s(results)
	return results[repeats/2] // Возвращаем медиану
}
\end{lstlisting}

Варьируя размер массива (\texttt{sizeBytes}), можно прицельно тестировать разные уровни кэша.
\textbf{Результаты для Apple M1:}
\begin{itemize}
    \item L1 Cache (размер теста 64 КБ): \textbf{0.9 ns}
    \item L2 Cache (размер теста 2 МБ): \textbf{3.5 ns}
    \item RAM (размер теста 128 МБ): \textbf{72 ns}
\end{itemize}

\paragraph{Пропускная способность ОЗУ.} Тест измеряет скорость, с которой процессор может читать или записывать большие объемы данных в оперативную память. Для этого создается большой слайс, и в цикле происходит его чтение или запись. 

\begin{lstlisting}[language=Go, caption={Реализация теста пропускной способности чтения}]
const dataSize = 1024 * 1024 * 64 // 64 MB

func ReadBandwidthBenchmark() (float64, string) {
	data := make([]byte, dataSize)
	// Заполняем, чтобы избежать оптимизации "нулевой" памяти
	for i := range data { data[i] = byte(i) }

	start := time.Now()
	var sink byte
	for _, v := range data {
		sink += v // Читаем каждый байт
	}
	elapsed := time.Since(start)
	
	mbps := float64(dataSize) / elapsed.Seconds() / (1024 * 1024)
	return mbps, "MB/s"
}
\end{lstlisting}

\textbf{Результаты для Apple M1 (LPDDR4X-4266):}
\begin{itemize}
    \item Пропускная способность чтения: \textbf{24.8 ГБ/с}
    \item Пропускная способность записи: \textbf{21.5 ГБ/с}
\end{itemize}

\subsubsection{Тесты механизмов параллелизма}

\paragraph{Производительность примитивов синхронизации.} Сравнивается производительность мьютекса (\texttt{sync.Mutex}) и атомарных операций (\texttt{atomic.AddInt64}) в условиях высокой конкуренции. Несколько горутин в цикле пытаются инкрементировать один общий счетчик.

\begin{lstlisting}[language=Go, caption={Реализация теста мьютекса}]
const syncOps = 1000000

func MutexBenchmark() (float64, string) {
	var wg sync.WaitGroup
	var mu sync.Mutex
	var counter int64
	numWorkers := runtime.NumCPU()
	wg.Add(numWorkers)

	start := time.Now()
	for i := 0; i < numWorkers; i++ {
		go func() {
			defer wg.Done()
			for j := 0; j < syncOps/numWorkers; j++ {
				mu.Lock()
				counter++
				mu.Unlock()
			}
		}()
	}
	wg.Wait()
	elapsed := time.Since(start)

	return float64(elapsed.Nanoseconds()) / syncOps, "ns/op"
}
\end{lstlisting}

\textbf{Результаты:}
\begin{itemize}
    \item Mutex Lock/Unlock: \textbf{25-40 ns/op}
    \item Atomic Increment: \textbf{2-5 ns/op}
\end{itemize}
Результаты показывают, что для простых операций, таких как инкремент счетчика, использование атомарных операций значительно эффективнее.
